% R3 pages of the R1-R3 exhibit book (block-r3-energy-emissions-correlates).
% Every number below is taken from the block's tables/*.csv, INDEX.md or the reviews' re-check sections.
% Development data -- obscured values, fictional firm panel. Not BPS findings.

\clearpage
\section{R3 --- Which firms cut energy use and CO2, and what goes with it}

{\small\textbf{Deasy's row R3 (PIC Toar), verbatim.}}
\begin{quote}\small
Request: ``Siapa yang reduce energy use and emisi? Lebih ke determinan motivasi apa yg membuat firms respond to climate
commitment.. Tambahan: Siapa yang impor lebih banyak (EGs) environment goods list (termasuk solar panel, dll)''.\\
Strategy: ``FE''.\quad Notes: ``Bukan kausal, tp mengetahui firms mana yg respond thd climate and motivasinya apa''.
\end{quote}

{\small\textbf{Every item of the row and where it is answered} (from the block's \texttt{TRACKER.md}).}\par\smallskip
{\footnotesize
\begin{tabularx}{\linewidth}{@{}>{\raggedright\arraybackslash}p{1.05cm}X>{\raggedright\arraybackslash}p{3.1cm}>{\raggedright\arraybackslash}p{6.3cm}@{}}
\toprule
ID & Deasy's item (short) & Exhibits & Status \\
\midrule
R3.Q1 & Who cuts energy use & R3-01 to 04, 06, 07, 09 & done \\
R3.Q2 & \ldots\ and emissions (CO2) & R3-01 to 04, 06, 07, 09 & done (total CO2 in every file) \\
R3.Q3 & What makes firms respond to the climate commitment & R3-02 to 07, 09 & partial: traits, place and stand-ins; motives are never asked \\
R3.Q3a & \quad size & R3-02, 04 & done \\
R3.Q3b & \quad export status & R3-02, 04 & done \\
R3.Q3c & \quad import status & R3-02, 04 & done, with caveat: the importer gap rests on changes starting in 2017 \\
R3.Q3d & \quad foreign ownership (10\% or more) & R3-02, 04 & done (2020 status not used) \\
R3.Q3e & \quad productivity & R3-07 & done \\
R3.Q3f & \quad R\&D, innovation & R3-05 & done (only the years that ask) \\
R3.Q3g & \quad environmental permits, waste treatment & R3-05 & done (asked in 2020--22 only) \\
R3.Q3h & \quad sector & R3-03, 04 & done \\
R3.Q3i & \quad province & R3-09, 04 & done (codes only, no names) \\
R3.Q3m & \quad making own electricity & R3-04, 06 & done \\
R3.Q3n & \quad starting energy level (high users fall back) & R3-01, 04 & done \\
R3.Q3o & \quad direct motives (climate pledge, buyers, parent firm) & --- & unavailable: no survey question asks \\
R3.Q4 & Who imports more environmental goods (solar panels etc.) & R3-08 & unavailable: no item-by-item import records \\
R3.S1 & Strategy ``FE'' & R3-04 & done \\
R3.N1 & ``Bukan kausal'': which firms respond, not cause & all R3 files & done (every file says association, not cause) \\
\bottomrule
\end{tabularx}}

\medskip
{\small\textbf{What this block says on the development data}}
\begin{itemize}\small
\item Each year roughly half of the firms record less energy and CO2 than the year before (36--72\% across years and
  measures). The jumps between years come mostly from masking and from which firms are observed, so from R3-02 on each change is compared with the typical
  change of the same year.
\item Per unit of value added, firm size, exporting, importing and foreign ownership barely matter: every group sits
  within 2 points of the typical firm (R3-02). Sectors do differ (R3-03); provinces do not (R3-09).
\item The two strongest traits, making own electricity and importing, rest on the changes from 2017 to 2018, when the
  recording of own-electricity firms' fuel shifted. Without them the gaps shrink or flip (R3-04, R3-06).
\item Permits, waste treatment, R\&D and innovation show no clear gap per unit of value added; gaps in totals come from
  big firms falling back after a high year (R3-05). Reducers look more productive mostly because of 2017 and because
  value added sits in both numbers (R3-07).
\item Imports of environmental goods cannot be measured: the file has one yearly total of imported raw materials per
  firm, with no list of items (R3-08).
\item Everything is association on masked values with fictional firm links: the pages test the method, not the answer.
\end{itemize}

% ---------------------------------------------------------------------------------------------------------------
\expage{R3-01}{Share of firms whose energy use and CO2 fell, each year}{\blockIII/figures/01_reducer_shares_by_year.pdf}{\texttt{01\_reducer\_shares\_by\_year.do} $\to$ \texttt{figures/01\_reducer\_shares\_by\_year.png},\newline \texttt{tables/01\_reducer\_shares\_by\_year.csv}, \texttt{tables/01\_reducer\_firms.csv}}{0.56}
\expart{What it is}{For every firm seen in two years in a row, whether its energy use (gigajoules) and its CO2 (tonnes)
were lower than the year before, in total and per unit of value added (output minus materials, energy and services
bought, current prices). Each point is the share of such firms in one year, 2011--2022: 20,169 firm pairs in all. The
file also builds the firm-level ``reducer'' flag that R3-02 to R3-07 use.}
\expart{How to read it}{Top panel energy, bottom panel CO2. Solid line with black dots = totals; dashed line with hollow
squares = per unit of value added. At 50\%, as many firms fell as rose. The gray band is 2016 (no annual file), so 2017
has no one-year change: the 2,049 firms seen in 2015 and 2017 are left out. The number under each year is the firm
pairs compared.}
\expart{What it shows on this data}{Before 2016 about half of the firms cut energy use each year (43.4--50.3\% in total,
50.1--56.0\% per unit of value added). After it the share swings: 63.0\% in 2018, 69.6\% in 2019, 43.6\% in 2022. These
swings mostly come from masking and from which firms are observed, not from behaviour. So the firm flag compares a
firm's changes with the typical (median) change of the same years: 1,865 of the 3,709 firms with at least two changes
(50.3\%) are reducers in energy per unit of value added, against 65.0\% by the raw count.}
\expart{How it helps Deasy's request}{It answers the first half of ``Siapa yang reduce energy use and emisi?'': how many
firms cut, each year, for energy and CO2, in total and per unit of value added. It does not say who they are; that
starts in R3-02, using the flag and the year-by-year comparison built here.}
\items{R3.Q1, R3.Q2, R3.N1}
\expart{Keep in mind}{It counts falls, not their size (counting only falls above 5\% lowers any year by at most 2.4
points). Price rises alone push the per-value-added share up. Firm links are fictional, so a ``firm pair'' may join two
real firms.}

% ---------------------------------------------------------------------------------------------------------------
\expage{R3-02}{Who cuts: size, exporters, importers, foreign owners}{\blockIII/figures/02_who_reduces_by_group.pdf}{\texttt{02\_who\_reduces\_by\_group.do} $\to$ \texttt{figures/02\_who\_reduces\_by\_group.png},\newline \texttt{tables/02\_who\_reduces\_by\_group.csv}, \texttt{tables/02\_who\_reduces\_firm\_flag.csv}}{0.50}
\expart{What it is}{The 20,169 one-year changes from R3-01, grouped by the firm's status at the start of each change:
size third (by workers, within the year), exporter, importer, and foreign-owned (10\% or more). For each group, the share
of its changes in which energy or CO2 fell more, or rose less, than the typical (median) firm's change that year. Four
panels: energy and CO2 per unit of value added, and in total.}
\expart{How to read it}{Dot = share; line = 95\% interval; dashed line at 50\% = same as the typical firm. Where the
interval crosses 50\%, the group is not clearly different. The number in brackets is the group's changes. ``Not
recorded'' rows are blank answers, never counted as non-exporters or domestic. Foreign status measured in 2020 is not
used: that year the foreign-owned share is more than twice the usual.}
\expart{What it shows on this data}{Per unit of value added every group sits within 2 points of 50\%; only non-exporters
(51.1\%, interval 50.2--52.0) and firms with no export answer (48.9\%) differ clearly, and only slightly. Exporters 49.3\%
against non-exporters 51.1\%, importers 50.3\% against 49.8\%, foreign-owned 48.5\% against domestic 50.1\%. In totals,
importers (51.9\%), the largest third (51.6\%) and, for CO2, foreign-owned firms (52.6\%, interval 50.1--55.2) cut more
often. But big starters fall back: changes from the top fifth of starting energy beat the typical change 69\% of the
time, the bottom fifth 30\%. The importer result also flips: 46.2\% before 2016, 51.5\% after (per unit of value added).}
\expart{How it helps Deasy's request}{It answers ``Siapa yang reduce energy use and emisi?'' for size, trade and
ownership, and the ``determinan'' part only as who, not why. On this data the answer is ``almost no difference per unit
of value added''.}
\items{R3.Q1, R3.Q2, R3.Q3, R3.Q3a, R3.Q3b, R3.Q3c, R3.Q3d, R3.N1}
\expart{Keep in mind}{Each trait is looked at alone; R3-04 holds them together. Comparing within the year stops the 2017
jump in the importer share from creating a gap by itself. Association, not cause.}

% ---------------------------------------------------------------------------------------------------------------
\expage{R3-03}{Who cuts, by manufacturing sector}{\blockIII/figures/03_who_reduces_by_sector.pdf}{\texttt{03\_who\_reduces\_by\_sector.do} $\to$ \texttt{figures/03\_who\_reduces\_by\_sector.png},\newline \texttt{tables/03\_who\_reduces\_by\_sector.csv}, \texttt{tables/03\_who\_reduces\_by\_sector\_firm\_flag.csv}}{0.56}
\expart{What it is}{The same measure as R3-02 for each two-digit manufacturing sector (the code at the start of each
change), over the same 20,169 changes, sorted by panel (a). Sectors with fewer than 100 changes (refined petroleum and
coke, 27; machine repair, 48) are named under the plot but not drawn.}
\expart{How to read it}{Dot = share of the sector's changes that beat the typical firm's change that year; line = 95\%
interval; dashed line = 50\%. An interval that misses 50\% marks a sector that is clearly different. Each change counts
once, whatever the firm's size.}
\expart{What it shows on this data}{Sectors do differ: in panel (a) 9 of 22 intervals miss 50\%, against about 1 by
chance. Per unit of value added, textiles (53.8\%) and paper (54.8\%) beat the typical firm more often; furniture
(43.5\%), beverages (43.0\%), computers and electronics (44.2\%) and electrical equipment (44.6\%) less often. Cement,
glass and ceramics cut total energy (46.9\%) and CO2 (47.2\%) slightly less often. Split at 2016, furniture and textiles
hold; steel (45.3\% before, 59.2\% after) and motor vehicles (44.0\%, 55.4\%) flip.}
\expart{How it helps Deasy's request}{It answers ``Siapa yang reduce energy use and emisi?'' by sector, and the
``determinan'' part as which sectors, not why. It answers the sector item fully on this data.}
\items{R3.Q1, R3.Q2, R3.Q3, R3.Q3h, R3.N1}
\expart{Keep in mind}{Value added is in current prices, so sectors whose prices rose faster look better in (a) and (b).
Whether sector codes are consistent across years is not checked. Each sector's share of all energy is in the table
only: one masked 2021 cement row holds 38.4\% of all energy, so do not name a ``biggest user''. Association, not cause.}

% ---------------------------------------------------------------------------------------------------------------
\expage{R3-04}{Which traits go with bigger cuts, other traits held fixed}{\blockIII/figures/04_trait_associations_fe.pdf}{\texttt{04\_trait\_associations\_fe.do} $\to$ \texttt{figures/04\_trait\_associations\_fe.png},\newline \texttt{tables/04\_trait\_associations\_fe.csv}, \texttt{tables/04\_trait\_associations\_fe\_all\_terms.csv}}{0.50}
\expart{What it is}{Deasy's ``FE'' strategy. For 17,922 one-year changes of 5,176 firms, the change in energy and CO2 (per
unit of value added, and in total) against the firm's traits at the start, all at once: workers, exporting, importing,
foreign ownership (10\% or more), making own electricity. Fixed effects mean comparing only like with like: (A) firms in
the same year, sector and province (main); (B) the same firm in its years with and without the trait. Changes starting
in 2020 are left out (foreign status unusable).}
\expart{How to read it}{Marker = how much more the measure changed for firms with the trait than for similar firms
without it, in \%; below 0 = bigger cut. Black dot = (A), hollow triangle = (B). Line = 95\% interval; crossing 0 = no
clear gap. Gray bar = where (A) moves when any one pair of years is left out. The bottom row, a starting level twice as
high, is a separate model.}
\expart{What it shows on this data}{In (A), making own electricity goes with a 36\% bigger cut in energy per unit of
value added (interval $-$41 to $-$31) and importing with 9\% ($-$16 to $-$1); workers, exporting and foreign ownership
show no clear gap. Both results rest on changes starting in 2017: without them own electricity is $-$18\% and the
importer gap flips to +14\%. In 2017 own-electricity firms recorded about 31 times other firms' fuel (1.4 to 4 times in
other years), likely a recording change. A starting level twice as high goes with a 38\% bigger cut, a third to a half
of it automatic fall-back.}
\expart{How it helps Deasy's request}{It follows the Strategy ``FE'' and answers ``determinan motivasi apa yg membuat
firms respond'' in part: which traits go with cuts, each net of the others, not motives.}
\items{R3.Q3, R3.Q3a, R3.Q3b, R3.Q3c, R3.Q3d, R3.Q3h, R3.Q3i, R3.Q3m, R3.Q3n, R3.S1, R3.N1}
\expart{Keep in mind}{(B) rests only on firms whose trait changes (importer 1,781, own electricity 1,943 firms), and such
switches may be recording shifts. A third version (sector-by-year and province-by-year) is in the tables only and moves
(A) by at most 3.2 points. Association, not cause.}

% ---------------------------------------------------------------------------------------------------------------
\expage{R3-05}{Motivation stand-ins: permits, waste treatment, R\&D, innovation}{\blockIII/figures/05_motivation_proxies.pdf}{\texttt{05\_motivation\_proxies.do} $\to$ \texttt{figures/05\_motivation\_proxies.png},\newline \texttt{tables/05\_motivation\_proxies.csv}, \texttt{tables/05\_motivation\_proxies\_fe.csv}}{0.55}
\expart{What it is}{The survey never asks why firms cut energy. The closest stand-ins are the rules a firm faces
(environmental permit documents AMDAL, UKL-UPL, SPPL, and waste treatment, asked only in 2020--22) and its capacity to
change (R\&D, any innovation, process innovation, asked in some years only). Each one-year change is grouped by the
answer at its start and compared with the typical change of firms that answered the same question that year.}
\expart{How to read it}{Black dot = share of the group's changes that beat that typical change; line = 95\% interval;
dashed line = 50\%. Compare ``Yes'' with ``No''. Hollow gray = no answer recorded, never counted as ``No''; these firms
differ in other ways, so their distance from 50\% is not a result. Permits: each firm sits under the strictest document
it holds (AMDAL, then UKL-UPL, then SPPL). 2017--18 R\&D and innovation codes are provisional.}
\expart{What it shows on this data}{Per unit of value added no stand-in clearly differs from ``No'': AMDAL 49.6\% against
none 49.6\%, waste treatment 50.5\% against 49.8\%, R\&D 47.2\% against 50.2\%, any innovation 50.1\% against 49.9\%. In
totals, permit holders and innovators cut a little more often (UKL-UPL or SPPL 54.5\% against 47.4\% in energy use), but
they start with 2--7 times more energy and high starts fall back; with starting energy held fixed, every total-energy
gap disappears or reverses.}
\expart{How it helps Deasy's request}{It is the closest this data gets to ``determinan motivasi apa yg membuat firms
respond to climate commitment''. It answers in part: rules faced and capacity to change, not motives. Direct motives
(climate pledge, buyer pressure, parent policy) are not asked (R3.Q3o).}
\items{R3.Q3, R3.Q3f, R3.Q3g, R3.N1}
\expart{Keep in mind}{Few years and only answering firms: against all firms, innovators would look worse (42.9\% against
49.0\%), within answering firms there is no gap. Association, not cause.}

% ---------------------------------------------------------------------------------------------------------------
\expage{R3-06}{Own electricity: reducers against non-reducers}{\blockIII/figures/06_own_electricity_reducers.pdf}{\texttt{06\_own\_electricity\_reducers.do} $\to$ \texttt{figures/06\_own\_electricity\_reducers.png},\newline \texttt{tables/06\_own\_electricity\_reducers\_by\_year.csv}, \texttt{tables/06\_own\_electricity\_reducers\_changes.csv}}{0.60}
\expart{What it is}{Part A: the share of reducers and non-reducers (firm flag from R3-01) that make some of their own
electricity, each year; left with all changes, right without changes starting in 2017. Part B: the share of one-year
changes that beat the typical firm's change, by own electricity at the start and from start to end.}
\expart{How to read it}{A: black dots = reducers, hollow squares = non-reducers, gray band = no 2016 file. B: black
dots = all changes, triangles = without changes starting in 2017, line = 95\% interval, dashed line = 50\%; hollow gray =
status not recorded (2014--15 only), not a result.}
\expart{What it shows on this data}{Reducers make own electricity more often: 66.7\% against 36.0\% in 2017, 27.1\%
against 16.6\% in 2022, +10.1 points over all years. This is mostly the 2017 recording change: that year own-electricity
firms recorded 31 times other firms' fuel, so their 2017-to-2018 changes look like big cuts. Without changes starting in
2017 the gap falls to +2.4 points, and in B(a) own-electricity firms beat the typical change in 49.0\% of changes
against 50.1\%. Holding the starting level fixed gives +4.3 points (energy per unit of value added) or $-$6.0 (energy use): no
difference that holds up either way.}
\expart{How it helps Deasy's request}{It answers ``Siapa yang reduce energy use and emisi?'' for self-generation
and the self-generation item under ``determinan motivasi'' (R3.Q3m). Answered on this data: the gap is mostly recording.}
\items{R3.Q1, R3.Q2, R3.Q3, R3.Q3m, R3.N1}
\expart{Keep in mind}{Self-made power is not in energy use, but starting to make it does not fake cuts: starters beat
the typical energy-use change in 47.9\% of changes. Own-electricity firms start with 4.4 times the energy, so gaps in
(c)--(d) are partly automatic. Association, not cause.}

% ---------------------------------------------------------------------------------------------------------------
\expage{R3-07}{Productivity of reducers against non-reducers}{\blockIII/figures/07_productivity_reducers.pdf}{\texttt{07\_productivity\_reducers.do} $\to$ \texttt{figures/07\_productivity\_reducers.png},\newline \texttt{tables/07\_productivity\_reducers\_firms.csv}, \texttt{tables/07\_productivity\_reducers\_changes.csv}}{0.61}
\expart{What it is}{Productivity = value added beyond what the firm's workers and last year's capital predict
(Levinsohn--Petrin method), against firms of the same sector and year. A: average productivity of reducers against
non-reducers (flags from R3-01). B: one-year changes in four equal productivity groups at the start. Covered: 2,740 of
3,709 flagged firms, 7,115 of 20,169 changes (12 of 24 sectors, never 2010 or 2017).}
\expart{How to read it}{A: dot = \% difference in productivity, reducers against non-reducers; triangle = reducers found
without changes starting in 2017; line = 95\% interval; crossing 0 = no clear difference. B: read as R3-02, one row per
productivity group.}
\expart{What it shows on this data}{Reducers in energy per unit of value added look 51\% more productive (interval 37 to
67), but only 13\% (2 to 26) without changes starting in 2017; for energy use and CO2 the gap becomes $-$7\% and $-$8\%,
not clear. Value added is in both numbers, so a good year raises productivity and lowers energy per value added at once.
B: the least productive group beat the typical change in 56.7\% of changes, the most productive 42.6\%, because high
value added falls back; in energy use there is no pattern (48.7, 53.0, 47.2, 50.0\%).}
\expart{How it helps Deasy's request}{It answers the productivity item under ``determinan motivasi'' (R3.Q3e).
Answered on this data: reducers look more productive mostly because of 2017 and arithmetic; productive firms do not go
on to cut energy use more often.}
\items{R3.Q1, R3.Q2, R3.Q3, R3.Q3e, R3.N1}
\expart{Keep in mind}{The file's other productivity measure (the ``common'' one) exists for only 665 firm-years; with it the gap is
25.6\% ($-$4.1 to 64.4), too uncertain to read. Covered firms are smaller (median 38 workers against 57). Nominal
values, fictional firm links. Association, not cause.}

% ---------------------------------------------------------------------------------------------------------------
\expage{R3-08}{Imports of environmental goods: not available in this data}{}{\texttt{08\_environmental\_goods\_imports.do} $\to$ \texttt{tables/08\_environmental\_goods\_imports\_unavailable.csv} (no figure)}{0.50}
{\footnotesize
\begin{tabularx}{\linewidth}{@{}>{\raggedright\arraybackslash}p{3.6cm}X@{}}
\toprule
Field (from the CSV) & Content \\
\midrule
Deasy's request & ``Tambahan: Siapa yang impor lebih banyak (EGs) environment goods list (termasuk solar panel, dll)'' \\
Status & unavailable \\
Answer & Not available in this dataset. Each firm reports only one yearly total for the imported raw materials it used,
  with no list of what was bought, so solar panels and other environmental goods cannot be picked out. \\
What is missing & A record of what each firm imported, item by item (customs product codes such as the 6-digit HS code),
  with the amount paid, for each year. \\
What exists instead & One yearly total per firm of imported raw materials used (\texttt{import\_value\_nominal};
  \texttt{import\_flag} = 1 when above zero), with no product, country or date detail. \\
Why the total cannot stand in & It mixes every imported item and cannot be split by product; the survey field is
  ``imported raw material used'', so imported machines such as solar panels are probably not in it at all. \\
What would be needed & (1) Customs import records per firm, item by item, linkable to this survey; (2) an agreed list of
  environmental goods with product codes, e.g.\ the APEC list of 54 products; (3) a table translating product codes
  between versions (HS, Indonesia's BTKI). \\
Not an answer & R1-01, R3-02 and R3-04 use the same total only to mark importers of any raw material. \\
Search result & Candidates found: none. Firm-years 33,049, of which 12,206 with imports; years 2010--2015, 2017--2022. \\
\bottomrule
\end{tabularx}}
\expart{What it is}{A one-row table that records why the environmental-goods question cannot be answered from this file,
what exists instead, and what data would answer it. The do-file searches every import variable, including Indonesian
labels, and finds no item-level import record.}
\expart{How to read it}{Each line is one column of the CSV. ``Answer'' is the plain reply to Deasy; ``What would be
needed'' is the shopping list for a future request to BPS or customs.}
\expart{What it shows on this data}{No firm-by-product import data: the only import number is a yearly total of imported
raw materials, above zero in 12,206 of 33,049 firm-years. It says nothing about solar panels or other environmental
goods.}
\expart{How it helps Deasy's request}{It answers ``Siapa yang impor lebih banyak (EGs) environment goods list'' with
``not possible here'', says why, and names the data that would make it possible (customs records per firm plus an agreed
list, such as the 54 products listed by APEC, the Asia-Pacific Economic Cooperation forum).}
\items{R3.Q4}
\expart{Keep in mind}{The importer rows in R1-01, R3-02 and R3-04 are importers of any raw material, not of
environmental goods. The importer share also jumps partway through the period, likely a recording change.}

% ---------------------------------------------------------------------------------------------------------------
\expage{R3-09}{Who cuts, by province and island group}{\blockIII/figures/09_who_reduces_by_province.pdf}{\texttt{09\_who\_reduces\_by\_province.do} $\to$ \texttt{figures/09\_who\_reduces\_by\_province.png},\newline \texttt{tables/09\_who\_reduces\_by\_province.csv}, \texttt{tables/09\_who\_reduces\_by\_province\_firm\_flag.csv}}{0.56}
\expart{What it is}{The same measure as R3-02 and R3-03 for each province (the code at the start of each change; the
file has no province names), grouped by island: first digit 1--2 Sumatra, 3 Java, 5--9 other islands. Provinces with at
least 100 changes are drawn on their own (12); smaller ones are pooled within their island (Sumatra 6 provinces, 305
changes; other islands 16 provinces, 474 changes).}
\expart{How to read it}{Black dot = province; black diamond = whole island group; hollow circle = smaller provinces
pooled. Line = 95\% interval; dashed line = 50\%, the typical firm. Java holds 88\% of the 20,169 changes, so its rows have
the narrowest intervals.}
\expart{What it shows on this data}{Per unit of value added no province differs: in panel (a) 0 of 14 rows miss 50\%,
against about 1 by chance, and the island groups sit at 49.0\% (Sumatra), 50.1\% (Java) and 49.8\% (other islands). In
totals a few rows miss 50\%. The largest, province 14 (58.8\% in energy use, interval 50.9 to 66.7), is mostly start
size: its changes start at 11.3 times the typical firm's energy that year, and within the same fifth of starting energy
it drops to 45.9\%.}
\expart{How it helps Deasy's request}{It answers ``Siapa yang reduce energy use and emisi?'' by place, and the
``determinan'' part as where, not why. The province item is answered on this data, by code only.}
\items{R3.Q1, R3.Q2, R3.Q3, R3.Q3i, R3.N1}
\expart{Keep in mind}{Codes are as recorded each year, with no boundary changes applied (code 65 appears only from
2015). Each province's share of all energy is in the table only: one masked row can dominate (province 63 holds 38.5\%
of all energy on 51 changes). Association, not cause.}
